\section{A complete eight-state ES-to-paired-quartet family}
The received local comparison \cite[Section 5]{ESRHReceived}
maps one length-four ES primary algebra to one of the two
paired-quartet primary algebras. Here we construct the whole
finite flat family, including the other four ES states and
their final collision. The endpoint is the formal quartet
polynomial of RD13; no assertion that its roots are zeros
of the arithmetic zeta function is made. All original
labels, factors, cofactor corrections and receiver metrics
are specified below. The original ES marking and polynomial
provenance remain \cite{ES488,Tao}.

Fix $p>0$ and $\gamma>0$. The exact reciprocal identity is
$4/p=1/p+1/(2p/5)+1/(2p)$. The ES tuple is therefore
$(p;x,y,z_{\rm ES})=(p;p,2p/5,2p)$; $p=5$ gives the
integral example $(5;5,2,10)$. This is not claimed to be
in a hard-prime subdomain. Explicitly place its four
root slots in the order
$(\mathrm M,\mathrm N,\mathrm T,\mathrm E)=(p,2p/5,p,2p)$.
This declares the permutation from $(p,x,y,z_{\rm ES})$;
no historical label is silently replaced. For $0\le z\le1$
and $\tau=1-z$, set
\begin{equation}
\begin{aligned}
 \alpha&=p\tau+z(\tfrac12+i\gamma),&
 \beta&=\tfrac{2p}{5}\tau+z(\tfrac12-i\gamma),\\
 \chi&=2p\tau+z(\tfrac12-i\gamma),&
 S&=2\alpha+\beta+\chi=\tfrac{22p}{5}\tau+2z>0,\\
 A_z&=-S^{-1},&
 h_z(r)&=A_z(r-\alpha)^2(r-\beta)(r-\chi).
\end{aligned}\tag{EF1}
\end{equation}
The marked slots are now $(\alpha,\beta,\alpha,\chi)$.
Expanding gives precisely
\begin{equation}
\begin{aligned}
 h_z(r)&=A_zr^4+r^3+B_zr^2+C_zr+D_z,\\
 B_z&=A_z[\alpha^2+2\alpha(\beta+\chi)+\beta\chi],\\
 C_z&=-A_z[\alpha^2(\beta+\chi)+2\alpha\beta\chi],\qquad
 D_z=A_z\alpha^2\beta\chi,\\
 h_0(r)&=-\frac5{22p}(r-p)^2(r-2p/5)(r-2p),\qquad
 h_1(r)=-\tfrac12[(r-\tfrac12)^2+\gamma^2]^2.
\end{aligned}\tag{EF2}
\end{equation}
In particular the coefficient of $r^3$ is one at every
parameter, and the leading coefficient is retained.
The entire family lies on the discriminant because
the M/T root is double throughout.

For all computations keep the explicit change and inverse
\begin{equation}
 m=(\beta+\chi)/2=\tfrac{6p}{5}\tau+z(\tfrac12-i\gamma),\quad
 x=r-m,\quad r=x+m,\quad
 d=\tfrac{4p}{5}\tau,\quad k=\alpha-m=-\tfrac p5\tau+2i\gamma z,
 \quad\kappa=k^2-d^2.
 \tag{EF3}
\end{equation}
Then $h_z=A_z(x-k)^2(x^2-d^2)$. Neither $k-d$
nor $k+d$ vanishes on the parameter interval: their
real parts are respectively $-p\tau$ and $3p\tau/5$,
and their imaginary part is $2\gamma z$. Vanishing
would require $\tau=z=0$, which is impossible.
Thus $\kappa$ is a unit on an open complex neighborhood
of the entire interval. Over the ring of polynomial
functions of $z$ with $S\kappa$ inverted, the algebra
\begin{equation}
 \mathscr S_z=\mathbb C[z,(S\kappa)^{-1}][r,t]/
                  (h_z(r),t^2-h_z'(r))
 \tag{EF4}
\end{equation}
is free of rank eight with the original basis RD1.
Successive division by the two relations proves freeness;
all eight coefficient states persist at both endpoints.

\section{Complete primary projectors, their inverse and original return}
Put $a_0=k^2+d^2$ in this section only; this symbol is
not the original source coordinate $a$. The exact two
primary projectors are
\begin{equation}
 e_U=\frac{(x^2-d^2)(3k^2-d^2-2kx)}{\kappa^2},\qquad
 e_L=\frac{(x-k)^2(a_0+2kx)}{\kappa^2},\qquad
 e_U+e_L=1.
 \tag{EF5}
\end{equation}
The last equation is an identity of polynomials.
Modulo $(x-k)^2$, the first is one, because it is
the first-order inverse of $x^2-d^2$; modulo
$x^2-d^2$ it is zero. The second has the reverse
congruences, since
$(a_0-2kx)(a_0+2kx)\equiv\kappa^2$ modulo $x^2-d^2$.
The two factors are coprime because $\kappa$ is a unit.
Consequently $e_U^2=e_U$, $e_L^2=e_L$, $e_Ue_L=0$
in the exact quartic algebra, including $d=0$.
Their central extension gives
$\mathscr S_z=e_U\mathscr S_z\oplus e_L\mathscr S_z$.

Here are every coefficient and the inverse of that
decomposition. In increasing powers of $x$, the
columns of $e_U,e_U(x-k)$ and $e_L,e_Lx$ are
\begin{equation}
 P_U=\frac1{\kappa^2}\begin{pmatrix}
 d^2(d^2-3k^2)&kd^2\kappa\\2kd^2&-d^2\kappa\\
 3k^2-d^2&-k\kappa\\-2k&\kappa
 \end{pmatrix},\qquad
 P_L=\frac1{\kappa^2}\begin{pmatrix}
 k^2a_0&2k^3d^2\\-2kd^2&k^2(k^2-3d^2)\\
 d^2-3k^2&-2k^3\\2k&a_0
 \end{pmatrix}.
 \tag{EF6}
\end{equation}
Multiplication by $x$ is reduced only by the original
quartic relation
$x^4=2kx^3-\kappa x^2-2kd^2x+k^2d^2$;
expanding EF5 and using that relation gives EF6.
For a cubic coefficient column $c$ the inverse is
\begin{equation}
 [P_U\ P_L]^{-1}c=
 \begin{pmatrix}1&k&k^2&k^3\\0&1&2k&3k^2\\
 1&0&d^2&0\\0&1&0&d^2\end{pmatrix}c,
 \qquad\det[P_U\ P_L]=\kappa^{-2}.
 \tag{EF7}
\end{equation}
The first two rows take the value and derivative at $x=k$.
The last two take the remainder modulo $x^2-d^2$.
Applying these four functionals to EF6 gives $I_4$;
the displayed square matrices therefore are mutual
inverses. Their determinant follows either by expansion
or the same four rows, whose determinant is $\kappa^2$.

In original $r$ coefficients use the full binomial
matrix $C_{-m}$ of MR1, since $x=r-m$. The complete
eight-dimensional primary inclusion and inverse are
\begin{equation}
Q_{\mathrm{prim}}=
 \operatorname{diag}(C_{-m}[P_U\ P_L],C_{-m}[P_U\ P_L]),\qquad
Q_{\mathrm{prim}}^{-1}=
 \operatorname{diag}([P_U\ P_L]^{-1}C_m,[P_U\ P_L]^{-1}C_m).
 \tag{EF8}
\end{equation}
The column order is upper even, lower even, upper odd,
lower odd, with two columns in each group in the order
shown in EF6. Its source and target returns are the exact
$F_UQ_{\mathrm{prim}}$ and $F_WQ_{\mathrm{prim}}$ of
GD14, with $T_A\oplus T_A$ between them. Its marked
column is $Q_r^{-1}Q_{\mathrm{prim}}z_{\mathrm{prim}}$.
All inverse coefficients are EF7--8 and FC24, and
their products are the identity. Thus the changing
root center $m$ is not discarded in the physical return.

The upper and lower primary relations are respectively
\begin{equation}
\begin{aligned}
 U_z&=\mathbb C[\varepsilon,t]/(\varepsilon^2,t^2-2g_0\varepsilon),
 &\varepsilon&=x-k,&g_0&=A_z\kappa,\quad g_1=2A_zk,\\
 L_z&=\mathbb C[x,t]/(x^2-d^2,t^2-u x-v),
 &u&=2A_z(k^2+d^2),&v&=-4A_zkd^2.
\end{aligned}\tag{EF9}
\end{equation}
The formulas follow by reducing the full derivative
$h_z'$ modulo the indicated root factors. Both algebras
have dimension four, with bases $(1,\varepsilon,t,t\varepsilon)$
and $(1,x,t,tx)$. Since $g_0\ne0$, the first is the
length-four algebra with $\varepsilon=t^2/(2g_0)$,
$t^4=0$ and basis $(1,t,t^2,t^3)$.
At $z=1$, $d=0$, $k=2i\gamma$, $A_z=-1/2$,
and the lower relation becomes $t^2=4\gamma^2x$,
$x^2=0$. Thus the whole endpoint has two length-four
primary algebras. At $z=0$ the lower roots are
distinct and $h_0'$ is nonzero there, so that summand
has four distinct signed states. The upper summand is
already nonreduced at every parameter.

\section{The complete residue corrections along the family}
The upper root cofactor at $\varepsilon=0$ is
$A_z[(\varepsilon+k)^2-d^2]=g_0+g_1\varepsilon+A_z\varepsilon^2$.
Expansion of its reciprocal through the exact quotient
$\varepsilon^2=0$ gives
\begin{equation}
 \Lambda_U(t(b_0+b_1\varepsilon))
   =\frac{b_1}{A_z\kappa}-\frac{2k b_0}{A_z\kappa^2},\qquad
 \Lambda_U(1)=\Lambda_U(\varepsilon)=0.
 \tag{EF10}
\end{equation}
For the lower component the inverse of its full cofactor
$A_z(x-k)^2$ modulo $x^2-d^2$ is
$(a_0+2kx)/(A_z\kappa^2)$. Taking the linear coefficient
of its product with $b_0+b_1x$ proves
\begin{equation}
 \Lambda_L(t(b_0+b_1x))
   =\frac{a_0b_1+2kb_0}{A_z\kappa^2},\qquad
 B_L=\begin{pmatrix}0&J_L\\J_L&0\end{pmatrix},\quad
 J_L=\frac1{A_z\kappa^2}
       \begin{pmatrix}2k&a_0\\a_0&2kd^2\end{pmatrix},\quad
 \det B_L=\frac1{A_z^4\kappa^4}.
 \tag{EF11}
\end{equation}
Both formulas are the restrictions of the full RD2
residue under the exact projectors EF5. They detect
every nilpotent at $d=0$. All cofactor terms are present.

At the original ES upper root, $g_{0,E}=3p/22$ and
$g_{1,E}=1/11$. Along the interval choose the continuous
square root $c(z)^2=g_{0,E}/g_0(z)$ with $c(0)=1$;
it exists since the nonzero continuous function has a
continuous argument on an interval. The complete local
algebra map and residue correction are
\begin{equation}
\begin{aligned}
 \phi_z(\varepsilon_E)&=\varepsilon,&\phi_z(t_E)&=c(z)t,\\
 \Lambda_U(\phi_z f)&=\Lambda_E(w_E(z)f),&
 w_E(z)&=c(z)^3\left[1+
       \left(\frac2{3p}-\frac{2k}{\kappa}\right)\varepsilon_E\right],\\
 w_E(z)^{-1}&=c(z)^{-3}\left[1-
       \left(\frac2{3p}-\frac{2k}{\kappa}\right)\varepsilon_E\right].
\end{aligned}\tag{EF12}
\end{equation}
The two relations in EF9 are preserved because
$c^2g_0=g_{0,E}$; the inverse sends $t$ to $c^{-1}t_E$
and $\varepsilon$ to $\varepsilon_E$. Evaluate both
residue functionals on the full odd basis
$t_E,t_E\varepsilon_E$ using EF10. Their two values
agree exactly by the displayed coefficient; the even
basis has zero residue. This proves the identity on
every element. Its stated inverse uses
$\varepsilon_E^2=0$. At $z=1$ the bracket is
$1+(2/(3p)+i/\gamma)\varepsilon_E$, and
$c(1)^2=3p/(44\gamma^2)$, reproducing the full
received upper local correction with its original sign.
EF5--11 retain the four lower states absent from a
comparison of just those two local blocks.

\section{The four remaining signed branches and their exact inverse cost}
For $z<1$, $d>0$ and the lower root/sign states are
\begin{equation}
 x_\epsilon=\epsilon d,\quad r_\epsilon=m+\epsilon d,\quad
 t_{\epsilon,\sigma}=\sigma(\epsilon d-k)\sqrt{2A_z\epsilon d},
 \qquad\epsilon,\sigma\in\{1,-1\}.
 \tag{EF13}
\end{equation}
Choose $\sqrt{2A_zd}=i\sqrt{-2A_zd}$ and
$\sqrt{-2A_zd}>0$. Both are continuous on the interval.
Then $t^2=2A_z\epsilon d(\epsilon d-k)^2=h_z'(r_\epsilon)$,
and every one of the four values is nonzero.
The exact original source state is FS13,
\begin{equation}
 q_{\epsilon,\sigma}=
 \left(t^{-1},-r-it,A_zt^3+2rt+3it^2,
 7ir^2t+(B_z-17r+A_zr^2)t^2-13it^3-2A_zt^4\right),
 \quad(r,t)=(r_\epsilon,t_{\epsilon,\sigma}).
 \tag{EF14}
\end{equation}
In particular these are actual finite inverse states of
the original polynomial for every $z<1$, and not just
points of a coefficient diagram. The N slot is
$\epsilon=-1$ and the E slot is $\epsilon=1$, with
both factor signs $\sigma$ retained. The M/T primary
state has $t=0$ throughout and belongs to the completed
boundary, not to an additional finite inverse branch.

As $z\to1$, $r_\epsilon\to1/2-i\gamma$ and
$t_{\epsilon,\sigma}/\sqrt d\to
\sigma(-2i\gamma)\sqrt{-\epsilon}$, with the square
roots just specified. Multiplication of every coordinate
in EF14 by $\sqrt d$ gives
\begin{equation}
 \sqrt d\,q_{\epsilon,\sigma}
 \longrightarrow\frac{\mathbf e_1}
       {\sigma(-2i\gamma)\sqrt{-\epsilon}},\qquad
 \sqrt d\,\|\Phi q_{\epsilon,\sigma}\|_U
 \longrightarrow\frac{\|\Phi\mathbf e_1\|_U}{2\gamma},\quad
 \sqrt d\,\|\Psi q_{\epsilon,\sigma}\|_W
 \longrightarrow\frac{\|\Psi\mathbf e_1\|_W}{2\gamma}.
 \tag{EF15}
\end{equation}
All remaining terms tend to zero by their displayed
powers of $t$ and bounded $r,A_z,B_z$. Fixed linearity
and continuity of the original norms prove both limits.
Equal norm limits therefore do not identify the four
labelled and signed branches or their distinct vector phases.

In the exact lower basis $(1,x,t,tx)$ let
\begin{equation}
 X_d=\begin{pmatrix}0&d^2\\1&0\end{pmatrix},\quad
 H_d=uX_d+vI,\quad D_d=v^2-u^2d^2=-4A_z^2d^2\kappa^2,\quad
 \Theta_L=M_{2t^3}=2\begin{pmatrix}0&H_d^2\\H_d&0\end{pmatrix}.
 \tag{EF16}
\end{equation}
Its exact inverse for $z<1$ is
\begin{equation}
 \Theta_L^{-1}=\tfrac12\begin{pmatrix}0&H_d^{-1}\\H_d^{-2}&0\end{pmatrix},\quad
 H_d^{-1}=\frac{vI-uX_d}{D_d},\quad
 H_d^{-2}=\frac{(v^2+u^2d^2)I-2uvX_d}{D_d^2},\quad
 \det\Theta_L=-1024A_z^6d^6\kappa^6.
 \tag{EF17}
\end{equation}
The identity $X_d^2=d^2I$ proves every product and the
inverse formulas. The determinant follows from
$\det M_t=\det H_d=D_d$ in this dimension, so
$\det M_{2t^3}=2^4D_d^3$. The lower trace Gram is
\begin{equation}
 T_L=B_L\Theta_L=
 \operatorname{diag}\left(
 \begin{pmatrix}4&0\\0&4d^2\end{pmatrix},
 \begin{pmatrix}4v&4ud^2\\4ud^2&4vd^2\end{pmatrix}\right),
 \qquad\det T_L=-1024A_z^2d^6\kappa^2.
 \tag{EF18}
\end{equation}
Multiplication in EF9 gives traces $4,0,0,0$ on
$(1,x,t,tx)$; reducing their ten products proves the
matrix directly and also verifies the residue factorization.
The full eight-dimensional trace is singular at every
parameter because of the upper nilpotents. Only the
specified lower four-dimensional inverse is used here.

The exact limiting inverse coefficient, with no phase
removed, is
\begin{equation}
 K_L=\lim_{z\to1}d^2\Theta_L^{-1}
  =\tfrac12\begin{pmatrix}
       0&(4\gamma^2)^{-1}N_2\\
       (16\gamma^4)^{-1}(I-2iN_2/\gamma)&0
     \end{pmatrix},\qquad
 N_2=\begin{pmatrix}0&0\\1&0\end{pmatrix},\qquad\operatorname{rank}K_L=3.
 \tag{EF19}
\end{equation}
Insert $u\to4\gamma^2$, $k\to2i\gamma$,
$A_z\to-1/2$ into EF17 after multiplying by $d^2$.
The upper block has rank one and the lower block is
invertible; their different parity images prove rank three.
This calculation retains the derivative of the cofactor
in the nonzero $N_2$ term of the lower block.

Here are the actual original norms for these constants.
Let $\Gamma_{r,1}=C_{1/2}^*G_{U,x}C_{1/2}$,
$V_L=C_{-m}P_L$, and
\begin{equation}
 G_L(z)=V_L^*\Gamma_{r,1}V_L,\qquad
 \Gamma_L(z)=\operatorname{diag}(G_L(z),G_L(z)),\qquad
 \Gamma_* =\Gamma_L(1),\quad G_*=G_L(1).
 \tag{EF20}
\end{equation}
The inclusion into the original source is
$J_L(z)=F_U\operatorname{diag}(V_L,V_L)$.
Equations EF6--8 prove it is injective through the
endpoint, so $G_L(z)$ is positive and converges to
the positive $G_*$. This is its induced original norm,
not a replacement metric. The target return uses
$F_W$ and $G_{W,x}$ with the same exact formula.
The equation $(T_A\oplus T_A)J_L(z)=J_{L,W}(z)$
proves the conductor comparison on this moving summand.

Set $L_*=\Gamma_*^{-1}K_L^*\Gamma_*K_L$.
Its three positive eigenvalues are $c_1^2\ge c_2^2\ge c_3^2>0$,
specified without missing metric entries by
\begin{equation}
\begin{aligned}
 a_1&=\operatorname{tr}L_*,\quad
 a_2=\tfrac12(a_1^2-\operatorname{tr}L_*^2),\quad
 a_3=\tfrac16(a_1^3-3a_1\operatorname{tr}L_*^2+2\operatorname{tr}L_*^3),\\
 \{c_1^2,c_2^2,c_3^2\}&=
    \{\text{roots of }s^3-a_1s^2+a_2s-a_3\},\qquad
 \eta=8\gamma^2\sqrt{(G_*)_{11}(G_*^{-1})_{00}}.
\end{aligned}\tag{EF21}
\end{equation}
Similarity with the positive semidefinite Euclidean
square of $K_L$ proves positivity and exactly one zero
eigenvalue. The trace identities are the elementary
symmetric functions, proving the complete cubic.
The limiting forward map is
$8\gamma^2\left(\begin{smallmatrix}0&0\\N_2&0\end{smallmatrix}\right)$.
Its unique positive singular value in $\Gamma_*$ is
$\eta$, by the norm of this rank-one map: its image
vector has squared norm $(G_*)_{11}$ and its input
coordinate functional has squared norm $(G_*^{-1})_{00}$.
Continuity of the positive square roots of $\Gamma_L(z)$,
EF19, and this rank-one forward limit prove the entire
inverse spectrum and every exterior asymptote:
\begin{equation}
\begin{aligned}
 \sigma(\Theta_L^{-1})&\sim
       (c_1d^{-2},c_2d^{-2},c_3d^{-2},\eta^{-1}),\\
 \lim_{z\to1}d^{2\min(j,3)}\|\wedge^j\Theta_L^{-1}\|
       &=\prod_{\nu=1}^{\min(j,3)}c_\nu\,
           \eta^{-\max(j-3,0)}\quad(1\le j\le4),\\
 \|\wedge^4\Theta_L^{-1}\|
       &=\frac1{1024|A_z|^6d^6|\kappa|^6},\qquad
 \frac{c_1c_2c_3}{\eta}=\frac1{65536\gamma^{12}}.
\end{aligned}\tag{EF22}
\end{equation}
The three growing values follow from the rank-three
inverse limit, and the last value is the reciprocal
of the unique positive forward value. Exterior norms
are products of the largest singular values. At full
rank the source and target are the same retained
metric, so the norm is the absolute reciprocal
determinant EF17. Substituting the exact endpoint
$|A_z|=1/2$, $|\kappa|=4\gamma^2$ gives the final
constant and proves agreement with the spectral product.
These are all four remaining ES branches and their
full primary inverse cost in the original receiver.
They do not assert invertibility of the already
singular upper block or identify this auxiliary
family with an arithmetic family of zeta zeros.

The exact constants in EF21 admit a full quadratic
evaluation because the two retained parity metrics
in EF20 are equal. Define
$\rho=(G_*)_{11}(G_*^{-1})_{00}>0$, retaining $G_*$
as the full moment matrix in EF20. Then
\begin{equation}
 \{c_1,c_2,c_3\}=\left\{
 \frac{\sqrt\rho}{8\gamma^2},\quad
 \frac{\sqrt{1+\rho/\gamma^2}+\sqrt\rho/\gamma}{32\gamma^4},\quad
 \frac{\sqrt{1+\rho/\gamma^2}-\sqrt\rho/\gamma}{32\gamma^4}
 \right\},\qquad \eta=8\gamma^2\sqrt\rho.
 \tag{EF23}
\end{equation}
The braces specify a multiset, to be sorted as in EF21.
To prove it, the upper block in EF19 has the single
positive singular value $\sqrt\rho/(8\gamma^2)$.
For $V=I-2iN_2/\gamma$, its determinant is one and
$\operatorname{tr}(G_*^{-1}V^*G_*V)=2+4\rho/\gamma^2$:
the two terms linear in $N_2$ vanish because its trace
and the trace of its adjoint are zero. Thus the two
singular values of $V$ have product one and squared
sum $2+4\rho/\gamma^2$. Solving these two positive
equations gives $\sqrt{1+\rho/\gamma^2}\pm\sqrt\rho/\gamma$.
The unchanged factor $1/(32\gamma^4)$ in EF19 proves
the last two constants. Their product and the first
constant give $c_1c_2c_3=\sqrt\rho/(8192\gamma^{10})$,
which proves the final identity in EF22 again with
every metric factor present.

There is also an exact finite-parameter evaluation,
not only a limiting estimate. For $0\le z<1$, put
\begin{equation}
\begin{aligned}
 a_H&=\operatorname{tr}(G_L(z)^{-1}H_d^*G_L(z)H_d),&
 b_H&=|D_d|^2,\\
 a_{H^2}&=\operatorname{tr}(G_L(z)^{-1}(H_d^2)^*G_L(z)H_d^2),&
 b_{H^2}&=|D_d|^4,\\
 \lambda_{F,\pm}&=\tfrac12(a_F\pm\sqrt{a_F^2-4b_F})
       &&(F=H,H^2),\\
 \sigma(\Theta_L^{-1})&=
 \left\{\frac1{2\sqrt{\lambda_{H,+}}},
 \frac1{2\sqrt{\lambda_{H,-}}},
 \frac1{2\sqrt{\lambda_{H^2,+}}},
 \frac1{2\sqrt{\lambda_{H^2,-}}}\right\}.
\end{aligned}\tag{EF24}
\end{equation}
The positive singular squares of a two-dimensional
map $F$ have sum $a_F$ and product $|\det F|^2$,
as follows by conjugating its metric square to a
positive Hermitian matrix. Its characteristic
polynomial therefore proves the first three lines.
Both $H_d,H_d^2$ are invertible by EF16, so all
four values are positive. The last line follows
from the exact two orthogonal parity blocks of
EF17 with the identical retained $G_L(z)$.
Ordering these four explicitly evaluated values
$v_1\ge v_2\ge v_3\ge v_4$ gives the exact
finite-parameter exterior norm
$\|\wedge^j\Theta_L^{-1}\|=\prod_{\nu=1}^jv_\nu$
for every $1\le j\le4$. Every original center,
cofactor, Gamma mass and inverse moment enters
through the displayed $A_z,k,d,V_L,G_L(z)$.

Finally the complete eight-state residue and trace
are evaluated in the full primary order specified
after EF8. Define
$J_U=\left(\begin{smallmatrix}-2k/(A_z\kappa^2)&1/(A_z\kappa)\\
1/(A_z\kappa)&0\end{smallmatrix}\right)$.
The complete matrices and ranks are
\begin{equation}
\begin{aligned}
 B_{\mathrm{prim}}&=\begin{pmatrix}
 0&\operatorname{diag}(J_U,J_L)\\
 \operatorname{diag}(J_U,J_L)&0\end{pmatrix},
 &\det B_{\mathrm{prim}}&=\frac1{A_z^8\kappa^8},\\
 T_{\mathrm{prim}}&=\operatorname{diag}\left(
 4,0,4,4d^2,0,0,
 \begin{pmatrix}4v&4ud^2\\4ud^2&4vd^2\end{pmatrix}\right),
 &\operatorname{rank}T_{\mathrm{prim}}&=
 \begin{cases}5,&0\le z<1,\\2,&z=1.\end{cases}
\end{aligned}\tag{EF25}
\end{equation}
The two residue cross blocks are exactly EF10--11;
their off-primary products vanish because $e_Ue_L=0$.
Each cross determinant is $-1/(A_z^2\kappa^2)$,
which proves the full determinant. Multiplication
by a scalar plus a nilpotent on the upper length-four
algebra has trace four times its scalar part, since
its matrix in powers of $t$ is triangular with that
scalar on every diagonal entry. Hence the upper
trace is $\operatorname{diag}(4,0,0,0)$. Combining
it with the full lower trace EF18 in the specified
primary order gives the displayed matrix. For $d>0$
the lower trace determinant is nonzero by EF18,
whereas at $d=0$ its rank is exactly one. This proves
both full ranks while the residue remains perfect.
Conjugation and bilinear congruence through EF8 and
the original receivers preserve all the stated ranks.
